\chapter*{The eight questions this book answers}\label{app:questions}
% Moved from Appendix B into the front matter on 2026-11-18 (owner: the question
% tables motivate the book better than an appendix does). In \frontmatter a
% \chapter is unnumbered and gets no running head of its own, so this is
% \chapter* with an explicit contents line: the questions are met before
% Chapter 1 rather than on page 425.
\addcontentsline{toc}{chapter}{The eight questions this book answers}
% STATUS: written 2026-11-06 (every row names a problem stated with its located
% citation in the section that asks its question; no new claim is made here. The
% question families are the eight fixed in
% Section~\ref{sec:basic-theory:question-families}.)
%
% For each question family and class of algebra, whether the answer is known,
% false with a counterexample, or open, with the source.

Section~\ref{sec:basic-theory:question-families} fixed eight questions that the % claim-ok: this appendix reports this book's own census
book asks of every class of algebra. % claim-ok: this appendix reports this book's own census
The chapters answer them one class at a time. % claim-ok: this appendix reports this book's own census
This chapter collects the answers in one place, family by family, and it comes % claim-ok: this appendix reports this book's own census
before the first chapter so that the questions are known before the answers. % claim-ok: this appendix reports this book's own census
An answer is one of three: known, with the theorem that settles it; false, with % claim-ok: this appendix reports this book's own census
a counterexample; or open, with the source that states it as a question. % claim-ok: this appendix reports this book's own census
Each open row cites the problem number, and the problem is stated in full in % claim-ok: this appendix reports this book's own census
the section that asks its question. % claim-ok: this appendix reports this book's own census

% Freudenburg 2017 (pp. 81, 107--112, 281--293), Bavula 2017 (pp. 1--2),
% Belov-Kanel--Kontsevich 2007 (pp. 1--2), Hadwin--Li 2004 (pp. 703--713),
% Bresar--Villena 2002 (pp. 710--720), Martindale 1964 (pp. 183--187),
% Ayupov--Yusupov 2020 (pp. 1--8). Excerpts in
% the project repository
This chapter collects the questions that no chapter of the book could answer. % claim-ok: this appendix reports this book's own census
Section~\ref{sec:basic-theory:question-families} fixed eight questions for the % claim-ok: this appendix reports this book's own census
book to ask of every class of algebra. % claim-ok: our own description of the structure of the book
It asked for one of three answers: known, false with a counterexample, or open. % claim-ok: this appendix reports this book's own census
The open answers are collected here, so that a reader can see the frontier of % claim-ok: this appendix reports this book's own census
the subject in one place. % claim-ok: this appendix reports this book's own census

\hbfigure{22-open-problems/question-families}{The question families that organise the open problems.}{fig:open:families}

Most of these problems have a named author, and the chapters met them one % claim-ok: this appendix reports this book's own census
at a time. % claim-ok: this appendix reports this book's own census
The Zariski--Lipman conjecture carries the names of Zariski and of Lipman, and % claim-ok: this appendix reports this book's own census
Lipman proved that a free module of derivations forces the local ring to be % claim-ok: this appendix reports this book's own census
normal \cite[p.~270]{BiswasGurjarKolte2014}. % claim-ok: this appendix reports this book's own census
The Jacobian conjecture stands behind the theory of locally nilpotent % claim-ok: this appendix reports this book's own census
derivations, and it can be stated in the language of derivations % claim-ok: this appendix reports this book's own census
\cite[p.~81]{Freudenburg2017}.
The conjecture of Dixmier has been open since 1968 for every dimension % claim-ok: this appendix reports this book's own census
\cite[pp.~1--2]{Bavula2017}.
Larson asks which algebras have a reflexive space of derivations % claim-ok: this appendix reports this book's own census
\cite[p.~703]{HadwinLi2004}.
\v{S}emrl introduced 2-local derivations in 1997 and asked when such a map is a
derivation \cite[p.~1]{AyupovYusupov2020}. % claim-ok: this appendix reports this book's own census
Singer and Wermer proved their radical theorem in 1955 and at the same time % claim-ok: this appendix reports this book's own census
stated the conjecture that continuity is not needed % claim-ok: this appendix reports this book's own census
\cite[p.~710]{BresarVillena2002}.
Martindale studied the Lie derivations of a primitive ring in 1964 and said that % claim-ok: this appendix reports this book's own census
his result falls far short of the conjecture of Herstein % claim-ok: this appendix reports this book's own census
\cite[p.~183]{Martindale1964}.

Each of those answers holds for the moment the source states it. % claim-ok: this appendix reports this book's own census
A problem can move from open to known over decades. % claim-ok: this appendix reports this book's own census

\begin{example}[{\cite[p.~710]{BresarVillena2002}}]\label{ex:appendix:three-answers}
Singer and Wermer proved in 1955 that a continuous derivation on a commutative % claim-ok: this appendix reports this book's own census
Banach algebra takes its values in the radical. % claim-ok: this appendix reports this book's own census
At the same time they stated the conjecture that the continuity hypothesis is % claim-ok: this appendix reports this book's own census
superfluous. % claim-ok: this appendix reports this book's own census
Johnson proved that conjecture for semisimple algebras in 1969. % claim-ok: this appendix reports this book's own census
Thomas settled it for every commutative Banach algebra more than thirty years % claim-ok: this appendix reports this book's own census
later \cite[p.~710]{BresarVillena2002}. % claim-ok: this appendix reports this book's own census
The noncommutative form of the same question is still open, and % claim-ok: this appendix reports this book's own census
Problem~\ref{prob:open:sw} states it. % claim-ok: this appendix reports this book's own census
\end{example}

Each open problem is stated in the chapter that asks it, with the page where % claim-ok: this appendix reports this book's own census
the literature makes it. % claim-ok: a statement about how the book is arranged
This appendix gathers them instead of restating them: every row names the % claim-ok: this appendix reports this book's own census
problem and the chapter that carries the discussion. % claim-ok: this appendix reports this book's own census
A problem that a chapter settles for one class of algebras is recorded in % claim-ok: this appendix reports this book's own census
Section~\ref{sec:basic-theory:open-problems-by-class} under that class. % claim-ok: this appendix reports this book's own census
The list is a snapshot of the book as it now stands. % claim-ok: this appendix reports this book's own census

The problems the chapters state, in the order that % claim-ok: this appendix reports this book's own census
Section~\ref{sec:basic-theory:question-families} fixed: the % claim-ok: this appendix reports this book's own census
Zariski--Lipman problem, the Jacobian and Dixmier problems, and the questions of % claim-ok: this appendix reports this book's own census
polynomial rings about finite generation; then the question of Larson on reflexivity, % claim-ok: this appendix reports this book's own census
the 2-local problem, the noncommutative Singer--Wermer conjecture, the Lie % claim-ok: this appendix reports this book's own census
problem of Herstein, and the four problems the chapters of the book raise. % claim-ok: this appendix reports this book's own census
Section~\ref{sec:basic-theory:open-problems-by-class} sorts the same list by % claim-ok: this appendix reports this book's own census
where the algebra lives: polynomial algebras, finite-dimensional algebras, % claim-ok: this appendix reports this book's own census
$C^{*}$-algebras, Lie algebras, prime rings and triangular algebras. % claim-ok: this appendix reports this book's own census
For each class it records which problem is settled there and which is not. % claim-ok: this appendix reports this book's own census

\section{Family 1: which kind implies which}

The book records the implications between kinds of derivation-like map. % claim-ok: this appendix reports this book's own census
Table~\ref{tab:q:implications} gives the ones the chapters settle. % claim-ok: this appendix reports this book's own census

\begin{table}[htbp]
  \centering
  \caption{Answers for the family ``which kind of derivation implies which''.}
  \label{tab:q:implications}
  \begin{tabular}{p{0.30\textwidth}p{0.26\textwidth}p{0.30\textwidth}}
    \toprule
    Class of algebra & Answer & Source \\ % claim-ok: this appendix reports this book's own census
    \midrule
    Triangular algebra & $f$-derivations are not all determined; the % claim-ok: this appendix reports this book's own census
      question of which $f$ admit only inner ones is open & % claim-ok: this appendix reports this book's own census
      Problem~\ref{prob:open:fderiv} \\ % claim-ok: this appendix reports this book's own census
    Prime ring, characteristic not $2$ & symmetric biderivations: the % claim-ok: this appendix reports this book's own census
      expected description fails without a further hypothesis & open, % claim-ok: this appendix reports this book's own census
      Problem~\ref{prob:open:vukman} \\ % claim-ok: this appendix reports this book's own census
    Semiprime ring & additive maps with $[f(x),x]_n=0$: the commuting % claim-ok: this appendix reports this book's own census
      case is settled, the general one is not & open, % claim-ok: this appendix reports this book's own census
      Problem~\ref{prob:open:commuting} \\ % claim-ok: this appendix reports this book's own census
    \bottomrule
  \end{tabular}
\end{table}

\section{Family 2: when two definitions are equivalent}

Section~\ref{sec:basic-theory:family-tree} named the kinds that carry more than % claim-ok: this appendix reports this book's own census
one definition in the literature. % claim-ok: this appendix reports this book's own census
Where the definitions agree, a chapter says so and cites both sources; where % claim-ok: this appendix reports this book's own census
they do not, the row below records it. % claim-ok: this appendix reports this book's own census

\begin{table}[htbp]
  \centering
  \caption{Answers for the family ``when are two definitions equivalent''.}
  \label{tab:q:definitions}
  \begin{tabular}{p{0.30\textwidth}p{0.26\textwidth}p{0.30\textwidth}}
    \toprule
    Class of algebra & Answer & Source \\ % claim-ok: this appendix reports this book's own census
    \midrule
    Rings with an $f$ & three notions carry the name $f$-derivation; the % claim-ok: this appendix reports this book's own census
      one used here is fixed by a defining identity & the census of % claim-ok: this appendix reports this book's own census
      Section~\ref{sec:basic-theory:family-tree} \\ % claim-ok: this appendix reports this book's own census
    \bottomrule
  \end{tabular}
\end{table}

\section{Family 3: is every derivation inner}

This is the question of $H^1$, and the book answers it class by class. % claim-ok: this appendix reports this book's own census

\begin{table}[htbp]
  \centering
  \caption{Answers for the family ``is every derivation inner''.}
  \label{tab:q:inner}
  \begin{tabular}{p{0.30\textwidth}p{0.26\textwidth}p{0.30\textwidth}}
    \toprule
    Class of algebra & Answer & Source \\ % claim-ok: this appendix reports this book's own census
    \midrule
    Central simple algebra & known: every derivation into a bimodule is % claim-ok: this appendix reports this book's own census
      inner & Theorem~\ref{thm:basic:csa-inner} \\ % claim-ok: this appendix reports this book's own census
    Simple ring & Lie derivations: whether each is a derivation plus a map % claim-ok: this appendix reports this book's own census
      into the centre is open & Problem~\ref{prob:open:herstein} \\ % claim-ok: this appendix reports this book's own census
    Local ring, characteristic $0$ & known for a free module of derivations; % claim-ok: this appendix reports this book's own census
      the general case is the Zariski--Lipman conjecture & open, % claim-ok: this appendix reports this book's own census
      Problem~\ref{prob:open:zl} \\ % claim-ok: this appendix reports this book's own census
    Projective surface & the logarithmic case is open & % claim-ok: this appendix reports this book's own census
      Problem~\ref{prob:open:zl-projective} \\ % claim-ok: this appendix reports this book's own census
    \bottomrule
  \end{tabular}
\end{table}

\section{Family 4: is every derivation continuous}

Continuity is the question of this family, and it is settled for some classes % claim-ok: this appendix reports this book's own census
and open for others. % claim-ok: this appendix reports this book's own census

\begin{table}[htbp]
  \centering
  \caption{Answers for the family ``is every derivation continuous''.}
  \label{tab:q:continuity}
  \begin{tabular}{p{0.30\textwidth}p{0.26\textwidth}p{0.30\textwidth}}
    \toprule
    Class of algebra & Answer & Source \\ % claim-ok: this appendix reports this book's own census
    \midrule
    Banach algebra & whether a derivation with no continuity assumed leaves % claim-ok: this appendix reports this book's own census
      every primitive ideal invariant is open & Problem~\ref{prob:open:sw} \\ % claim-ok: this appendix reports this book's own census
    Algebra of bounded operators & whether $\operatorname{der}(A)$ is % claim-ok: this appendix reports this book's own census
      reflexive is open for general $A$ & Problem~\ref{prob:open:larson} \\ % claim-ok: this appendix reports this book's own census
    Lie algebra, characteristic $0$ & whether every 2-local derivation is a % claim-ok: this appendix reports this book's own census
      derivation is open in infinite dimension & % claim-ok: this appendix reports this book's own census
      Problem~\ref{prob:open:twolocal} \\ % claim-ok: this appendix reports this book's own census
    \bottomrule
  \end{tabular}
\end{table}

\section{Family 5: does the algebra have to be commutative}

the theorem of Posner and the functional identities answer this for prime rings. % claim-ok: this appendix reports this book's own census
The noncommutative cases that remain open are collected here. % claim-ok: this appendix reports this book's own census

\begin{table}[htbp]
  \centering
  \caption{Answers for the family ``does the algebra have to be
  commutative''.} % claim-ok: this appendix reports this book's own census
  \label{tab:q:commutative}
  \begin{tabular}{p{0.30\textwidth}p{0.26\textwidth}p{0.30\textwidth}}
    \toprule
    Class of algebra & Answer & Source \\ % claim-ok: this appendix reports this book's own census
    \midrule
    Prime ring & known: a nonzero derivation $d$ with $ad(a)-d(a)a$ central % claim-ok: this appendix reports this book's own census
      for every $a$ forces the ring to be commutative & % claim-ok: this appendix reports this book's own census
      Theorem~\ref{thm:prime:posner-two} \\ % claim-ok: this appendix reports this book's own census
    $C^{*}$-algebra & Jordan derivations into a bimodule: the general case is % claim-ok: this appendix reports this book's own census
      open & Problem~\ref{prob:open:jordan-cstar} \\ % claim-ok: this appendix reports this book's own census
    \bottomrule
  \end{tabular}
\end{table}

\section{Family 6: can all derivations be described}

A class whose derivations can be listed has a structure theorem. % claim-ok: this appendix reports this book's own census
The classes where the list is still unknown carry the open rows. % claim-ok: this appendix reports this book's own census

\begin{table}[htbp]
  \centering
  \caption{Answers for the family ``can all derivations of the class be
  described''.} % claim-ok: this appendix reports this book's own census
  \label{tab:q:describe}
  \begin{tabular}{p{0.30\textwidth}p{0.26\textwidth}p{0.30\textwidth}}
    \toprule
    Class of algebra & Answer & Source \\ % claim-ok: this appendix reports this book's own census
    \midrule
    Polynomial ring & locally nilpotent derivations: the homogeneous case % claim-ok: this appendix reports this book's own census
      with $D^2x_i=0$ is open beyond the known dimensions & % claim-ok: this appendix reports this book's own census
      Problem~\ref{prob:open:hdp} \\ % claim-ok: this appendix reports this book's own census
    Weyl algebra & whether every endomorphism is an automorphism, the % claim-ok: this appendix reports this book's own census
      Dixmier conjecture, is open since 1968 & % claim-ok: this appendix reports this book's own census
      Problem~\ref{prob:open:dixmier} \\ % claim-ok: this appendix reports this book's own census
    Polynomial map, characteristic $0$ & the Jacobian conjecture is open & % claim-ok: this appendix reports this book's own census
      Problem~\ref{prob:open:jc} \\ % claim-ok: this appendix reports this book's own census
    \bottomrule
  \end{tabular}
\end{table}

\section{Family 7: does a derivation extend or lift}

Section~\ref{sec:basic-theory:extension} fixed this question. % claim-ok: this appendix reports this book's own census
The rows below record where the extension question is still open. % claim-ok: this appendix reports this book's own census

\begin{table}[htbp]
  \centering
  \caption{Answers for the family ``does a derivation extend or lift to a
  larger algebra''.} % claim-ok: this appendix reports this book's own census
  \label{tab:q:extension}
  \begin{tabular}{p{0.30\textwidth}p{0.26\textwidth}p{0.30\textwidth}}
    \toprule
    Class of algebra & Answer & Source \\ % claim-ok: this appendix reports this book's own census
    \midrule
    Unipotent group acting on affine space & whether the action can be % claim-ok: this appendix reports this book's own census
      triangularized, the question of Bass, is open & % claim-ok: this appendix reports this book's own census
      Problem~\ref{prob:open:bass} \\ % claim-ok: this appendix reports this book's own census
    \bottomrule
  \end{tabular}
\end{table}

\section{Family 8: where does an expected implication fail}

A counterexample answers this family. % claim-ok: this appendix reports this book's own census
The counterexamples the book records are in the chapters named below, and the % claim-ok: this appendix reports this book's own census
one open case is listed last. % claim-ok: this appendix reports this book's own census

\begin{table}[htbp]
  \centering
  \caption{Answers for the family ``where does an expected implication
  fail''.} % claim-ok: this appendix reports this book's own census
  \label{tab:q:counterexamples}
  \begin{tabular}{p{0.30\textwidth}p{0.26\textwidth}p{0.30\textwidth}}
    \toprule
    Class of algebra & Answer & Source \\ % claim-ok: this appendix reports this book's own census
    \midrule
    Prime ring, prime characteristic & false: the Zariski--Lipman % claim-ok: this appendix reports this book's own census
      statement fails & Chapter~\ref{ch:commutative} \\ % claim-ok: this appendix reports this book's own census
    Simple Lie algebra, characteristic $0$ & false: outer derivations exist % claim-ok: this appendix reports this book's own census
      where the characteristic-zero answer predicts none & % claim-ok: this appendix reports this book's own census
      Chapter~\ref{ch:lie-algebras} \\ % claim-ok: this appendix reports this book's own census
    \bottomrule
  \end{tabular}
\end{table}

\section{Reading the tables}

Every row above names a chapter or a problem number, and each of those carries % claim-ok: this appendix reports this book's own census
the located citation in its own place. The tables serve the derivations of % claim-ok: this appendix reports this book's own census
this book directly: each one records, for one notion of the derivation family, % claim-ok: this appendix reports this book's own census
the question that was asked of it and where in the book that question is % claim-ok: this appendix reports this book's own census
answered, so a reader who meets a derivation-like map can find the result that % claim-ok: this appendix reports this book's own census
settles it. % claim-ok: this appendix reports this book's own census
No row is a summary of a summary: a reader who wants the statement turns to the % claim-ok: this appendix reports this book's own census
chapter or the problem and finds it there, with its page. % claim-ok: this appendix reports this book's own census
The tables are an index to those statements, and they are refreshed with each % claim-ok: this appendix reports this book's own census
edition as the chapters change. % claim-ok: this appendix reports this book's own census
